\Askhsh{14536_SOLUTION}{1}
\begin{ylh}{1}
    \item [Ύλη από εικόνα]
\end{ylh}
\begin{wrapfigure}[12]{r}{8cm} % Adjust lines and width as needed for proper wrapping
    \centering
    \begin{minipage}[t]{3.8cm} % First triangle ABΓ
        \centering
        \begin{tikzpicture}[scale=0.7]
            % Define points for triangle ABΓ
            \tkzDefPoint(0,0){B}
            \tkzDefPoint(3,0){G} % Gamma
            % For Angle A = 48 degrees, base angles are (180-48)/2 = 66 degrees.
            % Apex A_y = (base_length/2) * tan(base_angle)
            \tkzDefPoint(1.5, {1.5 * tan(66)}){A}

            % Draw triangle
            \tkzDrawPolygon(A,B,G)

            % Mark equal sides AB = AΓ (double lines)
            \tkzMarkSegment[mark=||](A,B)
            \tkzMarkSegment[mark=||](A,G)

            % Label points
            \tkzLabelPoint[above](A){A}
            \tkzLabelPoint[below left](B){B}
            \tkzLabelPoint[below right](G){$\Gamma$}

            % Draw angle A
            \tkzMarkAngle[arc=ll,size=0.8cm,fill=maincolor!20](G,A,B)
            \tkzLabelAngle[pos=1.1](G,A,B){$48^\circ$}
        \end{tikzpicture}
    \end{minipage}
    \hfill % Space between minipages
    \begin{minipage}[t]{3.8cm} % Second triangle EΔZ
        \centering
        \begin{tikzpicture}[scale=0.7]
            % Define points for triangle EΔZ
            % Isosceles EΔ=EZ, Angle Z = 66 deg. This implies Angle Δ = 66 deg.
            % Apex E_y = (base_length/2) * tan(base_angle)
            \tkzDefPoint(0,0){Z}
            \tkzDefPoint(3,0){D} % Delta
            \tkzDefPoint(1.5, {1.5 * tan(66)}){E}

            % Draw triangle
            \tkzDrawPolygon(E,Z,D)

            % Mark equal sides EΔ = EZ (single lines)
            \tkzMarkSegment[mark=|](E,D)
            \tkzMarkSegment[mark=|](E,Z)

            % Label points
            \tkzLabelPoint[above](E){E}
            \tkzLabelPoint[below left](Z){Z}
            \tkzLabelPoint[below right](D){$\Delta$}

            % Draw angle Z
            \tkzMarkAngle[arc=ll,size=0.8cm,fill=maincolor!20](D,Z,E)
            \tkzLabelAngle[pos=1.1](D,Z,E){$66^\circ$}
        \end{tikzpicture}
    \end{minipage}
\end{wrapfigure}

ΛΥΣΗ

Σχεδιάζουμε δύο ισοσκελή τρίγωνα ΑΒΓ (ΑΒ = ΑΓ) και ΕΔΖ (ΕΔ = ΕΖ), τέτοια ώστε $\hat{A} = 48^\circ$, $\hat{Z} = 66^\circ$ και $AB = 3 \cdot E\Delta$.

Στο ισοσκελές τρίγωνο ΑΒΓ το άθροισμα των γωνιών του είναι $180^\circ$. Οπότε καθεμιά από τις γωνίες της βάσης του θα είναι ίση με $\frac{180^\circ - 48^\circ}{2} = \frac{132^\circ}{2} = 66^\circ$. Στο ισοσκελές τρίγωνο ΕΖΔ έχουμε ότι η γωνία $\hat{Z}$ της βάσης του είναι ίση με $66^\circ$, οπότε και η άλλη γωνία της βάσης θα είναι $66^\circ$. Δηλαδή $\hat{\Delta} = \hat{Z} = 66^\circ$. Τα τρίγωνα ΑΒΓ και ΕΖΔ έχουν τις δυο γωνίες στη βάση τους μία προς μία ίσες, οπότε είναι όμοια.

\begin{alist}
\item [\textbeta)]
\begin{rlist}
\item Στα όμοια τρίγωνα ομόλογες είναι οι πλευρές που είναι απέναντι από τις ίσες γωνίες. Οι λόγοι που σχηματίζονται είναι $\frac{AB}{EZ}$, $\frac{A\Gamma}{E\Delta}$ και $\frac{B\Gamma}{Z\Delta}$ οι οποίοι είναι ίσοι μεταξύ τους, αφού τα τρίγωνα είναι όμοια. Δηλαδή ισχύει ότι: $\frac{AB}{EZ} = \frac{A\Gamma}{E\Delta} = \frac{B\Gamma}{Z\Delta}$.
\item Ο λόγος των βάσεων είναι ο λόγος $\frac{B\Gamma}{Z\Delta}$ ο οποίος είναι ίσος με το λόγο $\frac{AB}{EZ}$.
\[ \frac{B\Gamma}{Z\Delta} = \frac{AB}{EZ} = \frac{3 \cdot E\Delta}{EZ} = \frac{3 \cdot EZ}{EZ} = 3 \]
Άρα ο ζητούμενος λόγος των βάσεων είναι ίσος με $3$.
\end{rlist}
\end{alist}